% =============================================================================
\chapter{Introduction}
\label{ch:introduction}
% =============================================================================

This guide documents every page, tab, modal and icon in the \emph{alphaN2N}
avatar-production platform. It is intended as an internal reference — a way to
onboard new team members, align engineering and product on existing behaviour,
and capture what the UI actually does (not what we think it does).

\begin{designnote}
The app is still evolving, and you will notice visible overlap between some
pages — for instance the same piece of data may be reachable from several
nav entries, or two pages may offer similar filters against slightly
different scopes. This is deliberate: we are exploring which surfaces best
fit the team's day-to-day workflow, and some ideas will be merged, renamed
or retired as that picture firms up. Treat overlap as a sign that the UX is
still settling, not a mistake.
\end{designnote}

% --------------------------------------------------------------------------
\section{Audience}
\label{sec:intro-audience}

\begin{itemize}
  \pagelement{Engineers}{joining a feature area and needing a map of what exists.}
  \pagelement{Product \& design}{auditing flows, spotting gaps, and checking whether the UI matches the PRDs.}
  \pagelement{Operations \& support}{cross-referencing what users see against backend behaviour.}
  \pagelement{QA}{building test plans from a canonical list of pages, tabs and modals.}
\end{itemize}

% --------------------------------------------------------------------------
\section{How to read this guide}
\label{sec:intro-how-to-read}

The guide follows the app's own navigation. Part I covers the shell and public
pages everyone sees. Parts II–VII walk the global sidebar groups in order
(Overview, Projects/Avatars, Content, Production, Review, Tools). Part VIII
covers the Admin area. Part X then takes a complete walk through the pipeline
workspace using the \texttt{x121 — adult content} pipeline as a worked example,
so the reader can see how every page re-scopes when entered via a pipeline
context.

Within each chapter:

\begin{itemize}
  \pagelement{Hero shot}{a full-page screenshot of the page at rest, taken on a seeded dev database.}
  \pagelement{Sections}{one sub-section per prominent UI region (header, filters, grid, side panels, etc.) with its own screenshot.}
  \pagelement{Tabs}{every tab on a tabbed page is screenshotted and described.}
  \pagelement{Modals}{every modal the page can open appears with a screenshot of its opened state (backdrop included). Shared modals also appear in Appendix~A.}
\end{itemize}

Two appendices round the guide out:

\begin{itemize}
  \pagelement{Appendix~A}{master index of every modal in the app, grouped by feature domain.}
  \pagelement{Appendix~B}{every icon used in the UI, grouped by category, with a sentence on what it means in context here (since the same icon can mean different things in different apps).}
\end{itemize}

% --------------------------------------------------------------------------
\section{Capture conventions}
\label{sec:intro-capture}

Screenshots are taken against the \emph{live} development database — no seed
script wipes or reshapes the data, so every shot reflects whatever projects,
avatars and generated assets currently exist in the running dev environment.
Capture specs are strictly read-only.

We use two capture passes:

\begin{enumerate}
  \item \textbf{Global pass} — logged in as super-admin, pipeline selector open, walking the global sidebar top-to-bottom.
  \item \textbf{Pipeline pass} — pipeline \texttt{x121 — adult content} selected, walking the pipeline-scoped sidebar. \texttt{x121} is the only pipeline deep-dived; other pipelines (e.g.\ \texttt{y122}) behave identically and are not captured separately.
\end{enumerate}

Public email templates and external-share-link variants are out of scope for
this edition.

Where a page appears in both global and pipeline contexts and renders the same
component, we document it once (in the global part) and note any pipeline-scope
differences inline.

\begin{designnote}
The goal of this guide is fidelity to what the UI actually does today, not
aspirational design. If you find something here that disagrees with the running
app, the app is right — raise a doc update.
\end{designnote}

\begin{futureidea}
Automate screenshot capture with Playwright so every chapter rebuilds from
seeded state on demand, and the PDF stays in sync as the UI evolves.
\end{futureidea}
